\pdfminorversion=7%
\documentclass[aspectratio=169,mathserif,notheorems]{beamer}%
%
\xdef\bookbaseDir{../../bookbase}%
\xdef\sharedDir{../../shared}%
\RequirePackage{\bookbaseDir/styles/slides}%
\RequirePackage{\sharedDir/styles/styles}%
%
%\gdef\searchSpace{\ensuremath{\mathbb{X}}}%
%
\protected\gdef\algobox#1#2{\resizebox{#1\paperwidth}{!}{%
\bgroup%
\fboxsep=0pt%
\fboxrule=2pt%
\mbox{\fcolorbox{white}{white}{\bgroup%
\fboxsep=2pt%
\fboxrule=2pt%
\mbox{\fcolorbox{hfuu-red!80}{hfuu-orange!60}{%
\mbox{#2}%
}}\egroup%
}}%
\egroup}}%
%
\title{Metaheuristic Optimization: State of the Art}%
%
%
\begin{document}%
\startPresentation%
%
\section{Introduction}%
%
\begin{frame}%
\frametitle{Introduction}%
\begin{itemize}%
\item Today, we will discuss the research frontier on \gls{metaheuristic} optimization.%
\item<2-> Therefore, we will first introduce what \gls{metaheuristic} optimization is and why we need.%
\item<3-> Then we will discuss the current \gls{SOTA} on several of the many, many, many domains where \glspl{metaheuristic} are applied.%
\end{itemize}%
\end{frame}%
%
\begin{frame}%
\frametitle{Metaheuristics}%
\begin{itemize}%
\item \Glspl{metaheuristic} are a class of algorithms used for \alert{optimization}.%
\end{itemize}%
\end{frame}%
%
\input{\sharedDir/slides/optimization_three_views.tex}%
%
\begin{frame}%
\frametitle{Views on Optimization, Part~1}%
\begin{itemize}%
\item Let's start with the first view: \alert{Superlatives}.%
\end{itemize}%
\end{frame}%
%
\input{\sharedDir/slides/optimization_superlatives.tex}%
%
\begin{frame}%
\frametitle{Views on Optimization, Part~2}%
\begin{itemize}%
\item Now let us consider the second view: \alert{Solving hard problems}.%
\item<2-> Later, we learn how to quantify \emph{hard}.%
\item<3-> For now, let's first think about how to define \emph{problem}?%
\end{itemize}%
\end{frame}%
%
\input{\sharedDir/slides/optimization_optimization_problem.tex}%
\input{\sharedDir/slides/optimization_business.tex}%
%
\begin{frame}%
\frametitle{Views on Optimization, Part~3}%
\begin{itemize}%
%
\item This brings us to the third view of optimization: We need to solve problems where equations are \inQuotes{not enough} and for which we cannot always find the best solution.%
%
\item<2-> Let's start with high school Maths, move to undergraduate computer science, and then approach this field of problems.%
\end{itemize}%
\end{frame}%
%
\definecolor{square-color}{HTML}{A03214}%
\definecolor{rectangle-color}{HTML}{0000A0}%
\definecolor{small-triangle-color}{HTML}{AA0078}%
\definecolor{large-triangle-color}{HTML}{008200}%
\definecolor{x-color}{HTML}{00788C}%
%
\protected\gdef\aRectangle{{\color{rectangle-color}{\ensuremath{A_{\text{\ding{122}}}}}}}%
\protected\gdef\aRectangleDiff{{\color{rectangle-color}{\ensuremath{A'_{\text{\ding{122}}}}}}}%
\protected\gdef\aRectangleMax{{\color{rectangle-color}{\ensuremath{\widehat{A_{\text{\ding{122}}}}}}}}%
\protected\gdef\aSquare{{\color{square-color}{\ensuremath{A_{\text{\ding{110}}}}}}}%
\protected\gdef\aSmallTriangle{{\color{small-triangle-color}{\ensuremath{A_{\text{\ding{116}}}}}}}%
\protected\gdef\aLargeTriangle{{\color{large-triangle-color}{\ensuremath{A_{\text{\ding{115}}}}}}}%
\protected\gdef\aX{{\color{x-color}{\ensuremath{x}}}}%
\protected\gdef\aXMax{{\color{x-color}{\ensuremath{\hat{x}}}}}%
\protected\gdef\aRectangleOfX{\ensuremath{\aRectangle(\aX)}}%
\protected\gdef\aRectangleOfXDiff{\ensuremath{\aRectangleDiff(\aX)}}%
%
\begin{frame}[t]%
\frametitle{A Problem that we can solve with high school Maths equations}%
\noindent\parbox{0.65\linewidth}{\noindent%
\begin{itemize}%
\item A \textcolor{rectangle-color}{rectangle} should be placed inside a \textcolor{square-color}{square} with side\\length~13cm as sketched.%
%
\item<2-> What area has the \emph{largest} such \textcolor{rectangle-color}{rectangle}?%
%
\item<3-> We can solve this by realizing that $\aRectangleOfX={\color{square-color}{\left(13\text{cm}^2\right)}}-2{\color{small-triangle-color}{\left(\tfrac{1}{2}\aX^2\right)}}-2{\color{large-triangle-color}{\left(\tfrac{1}{2}(13\text{cm}-\aX)^2\right)}}$\only<-3>{.}\uncover<4->{ and then finding the maximum by setting $0=\aRectangleOfXDiff=-4\aX+26\text{cm}$.}%
%
\item<5-> We get $\aXMax=6.5\text{cm}$\only<-5>{.}\uncover<6->{ and, hence, $\aRectangleMax=84.5\text{cm}^2$.}%
%
\item<7-> Maxima and minima are superlatives $=$ optima.%
%
\item<8-> Sometimes we can simply compute them and do not need fancy algorithms.%
\end{itemize}}%
%%
\strut\\[\paperheight]\strut%
%
\locateGraphic{1-2}{width=0.3\paperwidth}{\sharedDir/graphics/rectangle_problem/rectangle_problem_1}{0.67}{0.15}%
\locateGraphic{3-}{width=0.3\paperwidth}{\sharedDir/graphics/rectangle_problem/rectangle_problem_3}{0.67}{0.15}%
%
\end{frame}%
%

\definecolor{start-color}{HTML}{FF00C8}%
\definecolor{goal-color}{HTML}{00C0C0}%
\definecolor{path-red}{HTML}{FF0000}%
\definecolor{path-pink}{HTML}{FF00FF}%
\definecolor{path-green}{HTML}{008000}%
\definecolor{path-cyan}{HTML}{00C8C8}%
\definecolor{path-gray}{HTML}{646464}%
\protected\gdef\apath#1#2{{\color{path-#1}{#2}}}%
\protected\gdef\noapath{\protected\gdef\apath##1##2{##2}}%
%
\begin{frame}[t]%
\frametitle{Problems that algorithms can solve efficiently}%
%
\locate{2-}{\parbox{0.47\paperwidth}{\noindent%
\begin{itemize}%
%
\only<-9>{%
\item I am at the {\color{start-color}{starting point 合肥大学南二区北大门}}.%
%
\item<3-> I want to go to the {\color{goal-color}{goal:~食堂}}.
%
\item<4-> How do I get there the \alert{fastest}?%
%
\item<5-> We know the campus map.%
%
\item<6-> We want to compute the shortest path~(before actually walking it).%
%
\item<7-> We know all the intersections where I could make turns.%
%
\item<8-> There is no mathematical equation that we can use to solve this.%
}%
%
\item<9-> But we can use the \pgls{algoAstar}\cite{HNR1968AFBFTHDOMCP,P1997AAP}, which iteratively constructs a solution by adding a {\color{path-green}{heuristic}} to a {\color{path-red}{cost}} function.%
%
\only<-15>{%
\item<10-> For example, I could walk for \apath{red}{$27s$} to this intersection.%
%
\item<11-> Or for \apath{red}{$83s$} to that one.%
%
\item<12-> Or for \apath{red}{$195s$} to that one.%
}%
%
\only<-23>{%
\item<13-> We have 3~choices: (a)~$\apath{red}{27s}\only<17->{+\apath{green}{213s}}$, (b)~$\apath{red}{83s}\only<18->{+\apath{green}{221s}}$, (c)~$\apath{red}{195s}\only<19->{+\apath{green}{362s}}$.%
%
\item<14-> For each intersection, we can compute the airline distance~(\apath{green}{as the crow flies}) to the~{\color{goal-color}{食堂}}.%
%
\item<15-> The actual walking distance can never be shorter than that.%
%
\item<16-> \textbf{Adding }the two distances for each node gives us a good impression of how useful going there would be.%
%
\item<23-> (a)~it is\dots%
}%
%
\only<-64>{%
\item<24-> We continue doing this procedure until the first path to the goal is found.%
%
\item<56-> We found a first complete path from {\color{start-color}{start}} to {\color{goal-color}{goal}}.%
%
\item<57-> It has the total length~\apath{red}{$268s$}.%
%
\item<58-> Now we need to only check the remaining paths that might be shorter, those that are at least as long can be ignored.%
%
\item<64-> We confirmed the shortest possible path and it takes~$268s$.%
}%
%
\item<65-> In the worst case, the \pgls{algoAstar} needs to \inQuotes{look} at all intersections~(nodes) and streets~(edges) once.%
%
\item<66-> Usually, it is quite efficient to find shortest paths on maps.%
%
\item<67-> If we look for shortest paths that do not visit any node twice, where the edge distances are not negative, where the number of choices per node is limited, and where the graph and path are not too big/long -- then this algorithm is very efficient.%
%
\end{itemize}%
}}{0.51}{0.1}%075}%
%
\locateGraphic{}{width=0.545\paperwidth}{\sharedDir/graphics/hfuu_campus_map_astar/hfuu_campus_map_astar}{-0.02}{0.15}%
%
\locateGraphic{2-6}{width=0.545\paperwidth}{\sharedDir/graphics/hfuu_campus_map_astar/hfuu_campus_map_astar_01_start_goal}{-0.02}{0.15}%
\locateGraphic{8-}{width=0.55045\paperwidth}{\sharedDir/graphics/hfuu_campus_map_astar/hfuu_campus_map_astar_03_overlay}{-0.02421}{0.1449}%
%
\locateGraphic{7-9}{width=0.545\paperwidth}{\sharedDir/graphics/hfuu_campus_map_astar/hfuu_campus_map_astar_02_paths}{-0.02}{0.15}%
%
\locateGraphic{10}{width=0.545\paperwidth}{\sharedDir/graphics/hfuu_campus_map_astar/hfuu_campus_map_astar_04a}{-0.02}{0.15}%
\locateGraphic{11}{width=0.545\paperwidth}{\sharedDir/graphics/hfuu_campus_map_astar/hfuu_campus_map_astar_04b}{-0.02}{0.15}%
\locateGraphic{12-16}{width=0.545\paperwidth}{\sharedDir/graphics/hfuu_campus_map_astar/hfuu_campus_map_astar_04c}{-0.02}{0.15}%
\locateGraphic{17}{width=0.545\paperwidth}{\sharedDir/graphics/hfuu_campus_map_astar/hfuu_campus_map_astar_04d}{-0.02}{0.15}%
\locateGraphic{18}{width=0.545\paperwidth}{\sharedDir/graphics/hfuu_campus_map_astar/hfuu_campus_map_astar_04e}{-0.02}{0.15}%
\locateGraphic{19}{width=0.545\paperwidth}{\sharedDir/graphics/hfuu_campus_map_astar/hfuu_campus_map_astar_04f}{-0.02}{0.15}%
%
\locateGraphic{20}{width=0.545\paperwidth}{\sharedDir/graphics/hfuu_campus_map_astar/hfuu_campus_map_astar_04g}{-0.02}{0.15}%
\locateGraphic{21}{width=0.545\paperwidth}{\sharedDir/graphics/hfuu_campus_map_astar/hfuu_campus_map_astar_04h}{-0.02}{0.15}%
\locateGraphic{22}{width=0.545\paperwidth}{\sharedDir/graphics/hfuu_campus_map_astar/hfuu_campus_map_astar_04i}{-0.02}{0.15}%
%
\locateGraphic{23}{width=0.545\paperwidth}{\sharedDir/graphics/hfuu_campus_map_astar/hfuu_campus_map_astar_05}{-0.02}{0.15}%
%
\locateGraphic{24}{width=0.545\paperwidth}{\sharedDir/graphics/hfuu_campus_map_astar/hfuu_campus_map_astar_06a}{-0.02}{0.15}%
\locateGraphic{25}{width=0.545\paperwidth}{\sharedDir/graphics/hfuu_campus_map_astar/hfuu_campus_map_astar_06b}{-0.02}{0.15}%
%
\locateGraphic{26}{width=0.545\paperwidth}{\sharedDir/graphics/hfuu_campus_map_astar/hfuu_campus_map_astar_06c}{-0.02}{0.15}%
\locateGraphic{27}{width=0.545\paperwidth}{\sharedDir/graphics/hfuu_campus_map_astar/hfuu_campus_map_astar_06d}{-0.02}{0.15}%
%
\locateGraphic{28}{width=0.545\paperwidth}{\sharedDir/graphics/hfuu_campus_map_astar/hfuu_campus_map_astar_07a}{-0.02}{0.15}%
\locateGraphic{29}{width=0.545\paperwidth}{\sharedDir/graphics/hfuu_campus_map_astar/hfuu_campus_map_astar_07b}{-0.02}{0.15}%
%
\locateGraphic{30}{width=0.545\paperwidth}{\sharedDir/graphics/hfuu_campus_map_astar/hfuu_campus_map_astar_07c}{-0.02}{0.15}%
%
\locateGraphic{31}{width=0.545\paperwidth}{\sharedDir/graphics/hfuu_campus_map_astar/hfuu_campus_map_astar_08a}{-0.02}{0.15}%
\locateGraphic{32}{width=0.545\paperwidth}{\sharedDir/graphics/hfuu_campus_map_astar/hfuu_campus_map_astar_08b}{-0.02}{0.15}%
\locateGraphic{33}{width=0.545\paperwidth}{\sharedDir/graphics/hfuu_campus_map_astar/hfuu_campus_map_astar_09a}{-0.02}{0.15}%
\locateGraphic{34}{width=0.545\paperwidth}{\sharedDir/graphics/hfuu_campus_map_astar/hfuu_campus_map_astar_09b}{-0.02}{0.15}%
%
\locateGraphic{35}{width=0.545\paperwidth}{\sharedDir/graphics/hfuu_campus_map_astar/hfuu_campus_map_astar_10a}{-0.02}{0.15}%
\locateGraphic{36}{width=0.545\paperwidth}{\sharedDir/graphics/hfuu_campus_map_astar/hfuu_campus_map_astar_10b}{-0.02}{0.15}%
\locateGraphic{37}{width=0.545\paperwidth}{\sharedDir/graphics/hfuu_campus_map_astar/hfuu_campus_map_astar_10c}{-0.02}{0.15}%
%
\locateGraphic{38}{width=0.545\paperwidth}{\sharedDir/graphics/hfuu_campus_map_astar/hfuu_campus_map_astar_10d}{-0.02}{0.15}%
\locateGraphic{39}{width=0.545\paperwidth}{\sharedDir/graphics/hfuu_campus_map_astar/hfuu_campus_map_astar_11a}{-0.02}{0.15}%
%
\locateGraphic{40}{width=0.545\paperwidth}{\sharedDir/graphics/hfuu_campus_map_astar/hfuu_campus_map_astar_11b}{-0.02}{0.15}%
%
\locateGraphic{41}{width=0.545\paperwidth}{\sharedDir/graphics/hfuu_campus_map_astar/hfuu_campus_map_astar_11c}{-0.02}{0.15}%
%
\locateGraphic{42}{width=0.545\paperwidth}{\sharedDir/graphics/hfuu_campus_map_astar/hfuu_campus_map_astar_12a}{-0.02}{0.15}%
\locateGraphic{43}{width=0.545\paperwidth}{\sharedDir/graphics/hfuu_campus_map_astar/hfuu_campus_map_astar_12b}{-0.02}{0.15}%
\locateGraphic{44}{width=0.545\paperwidth}{\sharedDir/graphics/hfuu_campus_map_astar/hfuu_campus_map_astar_12c}{-0.02}{0.15}%
\locateGraphic{45}{width=0.545\paperwidth}{\sharedDir/graphics/hfuu_campus_map_astar/hfuu_campus_map_astar_12d}{-0.02}{0.15}%
%
\locateGraphic{46}{width=0.545\paperwidth}{\sharedDir/graphics/hfuu_campus_map_astar/hfuu_campus_map_astar_13a}{-0.02}{0.15}%
\locateGraphic{47}{width=0.545\paperwidth}{\sharedDir/graphics/hfuu_campus_map_astar/hfuu_campus_map_astar_13b}{-0.02}{0.15}%
%
\locateGraphic{48}{width=0.545\paperwidth}{\sharedDir/graphics/hfuu_campus_map_astar/hfuu_campus_map_astar_13c}{-0.02}{0.15}%
%
\locateGraphic{49}{width=0.545\paperwidth}{\sharedDir/graphics/hfuu_campus_map_astar/hfuu_campus_map_astar_14a}{-0.02}{0.15}%
\locateGraphic{50}{width=0.545\paperwidth}{\sharedDir/graphics/hfuu_campus_map_astar/hfuu_campus_map_astar_14b}{-0.02}{0.15}%
\locateGraphic{51}{width=0.545\paperwidth}{\sharedDir/graphics/hfuu_campus_map_astar/hfuu_campus_map_astar_14c}{-0.02}{0.15}%
\locateGraphic{52}{width=0.545\paperwidth}{\sharedDir/graphics/hfuu_campus_map_astar/hfuu_campus_map_astar_14d}{-0.02}{0.15}%
%
\locateGraphic{53}{width=0.545\paperwidth}{\sharedDir/graphics/hfuu_campus_map_astar/hfuu_campus_map_astar_15a}{-0.02}{0.15}%
\locateGraphic{54}{width=0.545\paperwidth}{\sharedDir/graphics/hfuu_campus_map_astar/hfuu_campus_map_astar_15b}{-0.02}{0.15}%
\locateGraphic{55}{width=0.545\paperwidth}{\sharedDir/graphics/hfuu_campus_map_astar/hfuu_campus_map_astar_15c}{-0.02}{0.15}%
%
\locateGraphic{56}{width=0.545\paperwidth}{\sharedDir/graphics/hfuu_campus_map_astar/hfuu_campus_map_astar_16}{-0.02}{0.15}%
%
\locateGraphic{57}{width=0.545\paperwidth}{\sharedDir/graphics/hfuu_campus_map_astar/hfuu_campus_map_astar_17}{-0.02}{0.15}%
\locateGraphic{58}{width=0.545\paperwidth}{\sharedDir/graphics/hfuu_campus_map_astar/hfuu_campus_map_astar_18}{-0.02}{0.15}%
\locateGraphic{59}{width=0.545\paperwidth}{\sharedDir/graphics/hfuu_campus_map_astar/hfuu_campus_map_astar_19}{-0.02}{0.15}%
\locateGraphic{60}{width=0.545\paperwidth}{\sharedDir/graphics/hfuu_campus_map_astar/hfuu_campus_map_astar_20}{-0.02}{0.15}%
\locateGraphic{61}{width=0.545\paperwidth}{\sharedDir/graphics/hfuu_campus_map_astar/hfuu_campus_map_astar_21}{-0.02}{0.15}%
\locateGraphic{62}{width=0.545\paperwidth}{\sharedDir/graphics/hfuu_campus_map_astar/hfuu_campus_map_astar_22}{-0.02}{0.15}%
\locateGraphic{63-}{width=0.545\paperwidth}{\sharedDir/graphics/hfuu_campus_map_astar/hfuu_campus_map_astar_23}{-0.02}{0.15}%
%
\end{frame}%
%
\begin{frame}%
\frametitle{Transition to Hard Problems}%
\begin{itemize}%
\item So we first looked at an optimization problem that can easily be solved to optimality using high school mathematics.%
%
\item<2-> Then we looked at a problem that can be solved to optimality using an algorithm that you should have learned during your Bachelor's studies.%
%
\item<3-> What about \emph{hard} problems?%
%
\item<4-> What about problems where traditional Maths and algorithms do not really help us much?%
%
\item<5-> Here we go.%
\end{itemize}%
\end{frame}%
%
%
\input{\sharedDir/slides/optimization_hard_problems.tex}%
%
\section{Metaheuristics}%
%
\input{\sharedDir/slides/optimization_metaheuristic_cycle.tex}%
%
\input{\sharedDir/slides/optimization_metaheuristic_components.tex}%
%
\gdef\sharedSlideTitle{Example:~RLS and the $(1+1)$~EA}%
\input{\sharedDir/slides/optimization_opoea_rls.tex}%
%
\xdef\sharedSlideTitle{Example:~Simulated Annealing}%
\input{\sharedDir/slides/optimization_sa.tex}%
%
%
\begin{frame}[t]%
\frametitle{Types of Metaheuristics}%
%
\begin{itemize}%
%
\item We can build many interesting algorithms just using such components~(and sometimes a bit more problem-specific information).
%
\item<2-> \glsFull{algoSLS}\only<3->{:%
\begin{itemize}%
\item Keep one solution~$x_c$, sample at least one new solution~$x_n$ from its neighborhood in each step, move to better solutions as discovered.%
%
\only<-5>{%
\item<4-> {\color{green!60!black}Advantage}: Fast (and sometimes simple)%
\item<5-> {\color{red!60!black}Disadvantage}: May get stuck at a local optimum%
}%
\end{itemize}%
}%
%
\item<6-> Population-based Methods\only<7->{:%
\begin{itemize}%
\item Maintain population~$S$ of $|S|=\mu\geq1$ solutions, sample set~$N$ of $|N|=\lambda\geq1$ new solutions in each step by creating modified copies of the elements in~$S$, then update $S$.%
\only<-9>{%
\item<8-> {\color{green!60!black}Advantage}: Less likely to get stuck at local optimum.%
\item<9-> {\color{red!60!black}Disadvantage}: Slow.%
}%
\end{itemize}%
}%
%
\item<10-> Model-based Methods\only<11->{:%
\begin{itemize}%
\item Maintain a stochastic model~$M$ (usually a multivariate probility distribution) of what good solutions could look like, sample set~$N$ of $|N|=\lambda\geq1$ new solutions based on~$M$ in each step, then update~$M$.%
\only<-13>{%
\item<12-> {\color{green!60!black}Advantage}: Sometimes really good performance.%
\item<13-> {\color{red!60!black}Disadvantage}: Complicated, sampling and model update may be slow.%
}%
\end{itemize}%
}%
%
\item<14-> Hybrid Algorithms\only<15->{:%
\begin{itemize}%
\item Usually combine model-based or population-based approach with \glsFull{algoLS}: Each new solution from model/population is refined with local search.%
\only<-17>{%
\item<16-> {\color{green!60!black}Advantage}: Fast~(because \glsShort{algoLS}), good global search~(because of population/model)%
\item<17-> {\color{red!60!black}Disadvantage}: usually few.%
}%
\end{itemize}%
}%
%
\end{itemize}%
%
\end{frame}%
%
%
\section{Discrete Optimization}%
\subsection{Black-Box Discrete Optimization Benchmarks}%
%
%
\begin{frame}%
\frametitle{Black-Box Discrete Optimization Benchmarks: Introduction}%
\begin{itemize}%
%
\item \alert{Problem}: Search space~\searchSpace\ corresponds to the set~$\{0,1\}^{\scale}$ of bit strings of length~\scale\ and \emph{nothing} is known about objective~function~$\objFun:\{0,1\}^n\mapsto\naturalNumbersZ$.%
%
\item<2-> Benchmarks here often have the goal to investigate problematic aspects of discrete problems in isolation\cite{ZWWW2026WFFAFTSIFGLS,DYHWSB2020BDOHWI}.%
%
\item<3-> The benchmarks are functions that scale with~\scale\ and are usually computable in~\bigOb{\scale} or~\bigOb{\scale^2}.%
%
\item<4-> Their optima are known~(but the algorithms we apply to them don't know them).%
%
\end{itemize}%
\end{frame}%
%
\begin{frame}[t]%
\frametitle{Black-Box Discrete Optimization Benchmarks: Benchmarks}%
\begin{itemize}%
%
\item \alert<-1>{\gls{optOneMax}}\cite{DJW2006UALBFRSHIBBO,M1992HGARWMAH}\only<-1>{:~maximize number of \texttt{1} bits, very easy, single optimum.}%
%
\only<-2>{\item<2->}\only<3->{,} \alert<-2>{\glsFull{optLinHarm}}\cite{DYHWSB2020BDOHWI,DJW2002OTAOTOPOEA}\only<-2>{:~like \gls{optOneMax}, but bits have linearly increasing weights.}%
%
\only<-3>{\item<3->}\only<4->{,} \alert<-3>{\glsFull{optBinInt}}\cite{TGG1998DCDATSSOP}\only<-3>{:~like \gls{optLinHarm}, but bits have exponentially increasing weights.}%
%
\only<-4>{\item<4->}\only<5->{,} \alert<-4>{\gls{optLeadingOnes}}\cite{R1997CPOEA,W1989TGAASPWRBAORTIB}\only<-4>{:~number of leading \texttt{1} bits; contribution of variable~$x_i$ depends on all variables~$x_j$ with~$j<i$.}%
%
\only<-5>{\item<5->}\only<6->{,} \alert<-5>{\gls{optTrap}}\cite{DJW2002OTAOTOPOEA,NB2003AAOTBOSEAOTF}\only<-5>{:~\gls{optOneMax}, but optimum and worse solution swapped -- extremely hard, local search needs exponential runtime.}%
%
\only<-6>{\item<6->}\only<7->{,} \alert<-6>{\gls{optTwoMax}}\cite{FQW2018ELDBOAWHTMO,VHGN2002FTTIMEAHSP}\only<-6>{:~two optima, a global and a local one, of same size, local search needs exponential runtime with 50\% probability.}%
%
\only<-7>{\item<7->}\only<8->{,} \alert<-7>{\gls{optJump}}\cite{DJW2002OTAOTOPOEA,FQW2018ELDBOAWHTMO}\only<-7>{:~like \gls{optOneMax}, but last $\omega$~bit flips around optimum are deceptive = lead away from optimum.}%
%
\only<-8>{\item<8->}\only<9->{,} \alert<-8>{\gls{optPlateau}}\cite{AD2018PRAFP}\only<-8>{:~like \gls{optJump}, but last $\omega$~bit fips around optimum are neutral = have same objective value.}%
%
\only<-9>{\item<9->}\only<10->{,} \alert<-9>{\glsFullpl{optIsing}}\cite{DYHWSB2020BDOHWI,F2004APUBFAMBAOTTDIM,FW2005TODIMMVR}\only<-9>{:~variables should take on same values as neighbors in a given tology~(ring, torus, \dots); many interactions between variables.}%
%
\only<-10>{\item<10->}\only<11->{,} \alert<-10>{\gls{optNQueens}}\cite{DYHWSB2020BDOHWI}\only<-10>{:~place $N$~queens on a chessboard so that they cannot beat each other, many interactions between variables.}%
%
\only<-11>{\item<11->}\only<12->{,} \alert<-11>{\gls{optWModel}}\cite{WW2018DFOCOPATTWBPFST,WNSRG2008ATMFMOERANFL}\only<-11>{:~allows you to tune and combine the above features.}%
%
\item<12-> If an algorithm performs well on all of these problems, then it can probably also solve real problems quite well\dots%
%
\end{itemize}%
\locateGraphic{1}{width=0.3\paperwidth}{\sharedDir/graphics/discrete_benchmarks/discrete_benchmarks_onemax}{0.67}{0.6}%
\locateGraphic{5}{width=0.3\paperwidth}{\sharedDir/graphics/discrete_benchmarks/discrete_benchmarks_onemax_trap}{0.67}{0.6}%
\locateGraphic{6}{width=0.3\paperwidth}{\sharedDir/graphics/discrete_benchmarks/discrete_benchmarks_twomax}{0.67}{0.6}%
\locateGraphic{7}{width=0.3\paperwidth}{\sharedDir/graphics/discrete_benchmarks/discrete_benchmarks_jump}{0.67}{0.6}%
\locateGraphic{8}{width=0.3\paperwidth}{\sharedDir/graphics/discrete_benchmarks/discrete_benchmarks_plateau}{0.67}{0.6}%
\end{frame}%
%
\begin{frame}%
\frametitle{Black-Box Discrete Optimization Benchmarks: Baselines}%
\begin{itemize}%
%
\item \glsFull{algoRLS}:~unary operator flips 1~bit%
%
\item<2-> \glsShort{algoEAopo}:~unary operator flips each bit with probability~$1/n$%
%
\item<3-> \glsFull{algoSGA}\cite{G1989GA,W2009GOATAA,BFM1997HOEC}:~traditional population-based \glsShort{algoGA}, now outdated%
%
\item<4-> \glsFull{algoTS}:~like \gls{algoRLS}, but forbit flipping of recently flipped bits to avoid cycles%
%
\end{itemize}%
\end{frame}%
%
%
\begin{frame}%
\frametitle{Black-Box Discrete Optimization Benchmarks: SOTA}%
\begin{itemize}%
%
\item The \alert<-1>{\glsShort{algoSAGA}} by \citeauthor{CPD2018TAMPARAOEA}~\bracketCite{CPD2018TAMPARAOEA,DD2018OSASAPCFT1LLGA}: linear runtime on \gls{optOneMax}.%
%
\item<2-> The \alert<-2>{\glsShort{algoG2P1GA}} by \citeauthor{S2012CSUBBA}~\bracketCite{S2012CSUBBA,S2017HCSUBBAIGA,CPD2018TAMPARAOEA,CPD2018ASPFTUOCIBBO}:~An algorithm using a binary operator and proven to be twice as fast as any unary-operator only black-box method on \gls{optOneMax}.%
%
\item<3-> \alert<-3>{\glsShort{algoFEAopo}} by Weise et al.~\bracketCite{WWLC2021FFAMOAIUBTOTOFV,WWLCL2023FFAOWBFGSCBE,ZWWW2026WFFAFTSIFGLS}, i.e., the \glsShort{algoEAopo} with \glsFull{algoFFA}:~the \emph{only} algorithm achieving polynomial runtime on \gls{optTrap}, \gls{optTwoMax}, and \gls{optJump}.\uncover<4->{ %
\alert{Invented here.}%
}%
%
\end{itemize}%
%
\end{frame}%
%
\subsection{Maximum Satisfiability Problem}%
%
\xdef\sharedSlideTitle{Maximum Satisfiability Problem: Introduction}%
\input{\sharedDir/slides/optimization_maxsat.tex}%
%
\begin{frame}%
\frametitle{Maximum Satisfiability Problem: Benchmarks}%
\begin{itemize}%
%
\item At so-called phase transition region with~$\frac{m}{n}\approx4.26$, the average instance hardness for \gls{algoSLS} algorithms is maximal\cite{HS2000SAORFROS,DNS2017TCAOEAORSkCF,DNS2015IRBFTOPOEOR3CFBOFDC}, so instances from this region are most interesting.%
%
\item<2-> \gls{SATLIB}\cite{HS2000SAORFROS,HS2001STSL} provides satisfiable benchmark instances from the phase transition region~(i.e., the hardest type of instances) for different scales~$n\in20 \cup \{25 i\; \forall i \in \intRange{2}{10}\}$.%
%
\end{itemize}%
\end{frame}%
%
%
\begin{frame}%
\frametitle{Maximum Satisfiability Problem: Baselines}%
\begin{itemize}%
\item As black-box baselines, we can use the algorithms already listed for discrete benchmarks.%
 %
\item<2-> A good problem-specific baseline is \gls{algoWalkSAT} by \textcite{SKC1994NSFILS}.%
%
\end{itemize}%
\end{frame}%
%
%
\begin{frame}%
\frametitle{Maximum Satisfiability Problem: SOTA}%
\begin{itemize}%
%
\item For smaller-scale instances we can use \glsFull{algoBB} algorithms like \textcite{LMMP2010RBLBIM}.%
%
\item<2-> Encoding-based methods like those of \textcite{PRB2018DPWEFSWM} can deal with even much larger instances.%
%
\item<3-> The \gls{SOTA} for incomplete solvers\cite{GKSS2008SS} comprises \glsShort{algoSLS} algorithms, often based on ideas dating back to \gls{algoWalkSAT}, such as \glsFull{algoAMLS} by \textcite{LH2012AMBLSFMS} or \gls{algoNuWLS} by \textcite{CCL2023NILSFWPMBNWT}.%
%
\end{itemize}%
\end{frame}%
%
%
\subsection{Low Autocorrelation Binary Sequences}%
%
\begin{frame}%
\frametitle{Low Autocorrelation Binary Sequences: Introduction}%
\begin{itemize}%
\item \glsFullpl{optLABS} are bit strings of length~\scale\ whose autocorrelation is as low as possible.%
%
\item<2-> This means that sub-sequences of any length~$1\leq k<\scale$ should ideally never repeat.%
%
\item<3-> The objective function~$\objFun:\{0,1\}^{\scale}\mapsto\naturalNumbersZ$ represents the energy of a bit string~\sespel\ and can be defined as~$\objFunb{\sespel}=\sum_{k=1}^{\scale-1}\left(\sum_{i=0}^{\scale-1-k} \arrayIndex{s}{i} \cdot \arrayIndex{s}{i+k}\right)^2$ with~$\arrayIndex{s}{j}=2\arrayIndex{\sespel}{j}-1$.%
%
\item<4-> Such sequences are needed for communication engineering\cite{MDW2015NESFLLABS,BBB2017LABSOIMFARPTAT}, cryptography\cite{MDW2015NESFLLABS}, as well as modulation pulses in radar and sonar ranging\cite{PM2016LABS}.%
%
\item<5-> The \glsShort{optLABS} is \npHard\cite{T2021CNAOTELOTLABSP}.%
%
\end{itemize}%
\end{frame}%
%
\begin{frame}%
\frametitle{Low Autocorrelation Binary Sequences: Benchmarks}%
\begin{itemize}%
\item The problem itself is also the benchmark.%
\item<2-> So we only got the objective function~\objFun\ for different values of~\scale\dots%
\end{itemize}%
\end{frame}%
%
%
\begin{frame}%
\frametitle{Low Autocorrelation Binary Sequences: Baselines}%
\begin{itemize}%
\item The algorithms listed under discrete optimization benchmarks.%
%
\item<2-> For small scales~\scale\ the \gls{algoGA} by \textcite{MZB1998ESFLABS} or the exhaustive enumeration by \textcite{M1996ESFLABS}.%
%
\item<3-> See also the (outdated) online repository \citeurl{BBB2016AGAFSASOTLP} by \textcite{BBB2016AGAFSASOTLP}.%
\end{itemize}%
\end{frame}%
%
%
\begin{frame}%
\frametitle{Low Autocorrelation Binary Sequences: SOTA}%
\begin{itemize}%
%
\item the algorithm by \textcite{D2022NCOBSWHMF}\uncover<2->{,}%
%
\item<2-> the \glsFull{algoMA} by \textcite{MDW2015NESFLLABS}\uncover<3->{,}%
%
\item<3-> the xLastovka-family of methods by \citeauthor{BB2022CSOLSSBSWHMF}~\bracketCite{BB2022CSOLSSBSWHMF,BB2018AHAFALABSPWOLAHMF}\uncover<4->{, and}%
%
\item<4-> the \glsShort{algoSLS} by \textcite{DBN2020EGOLABS}~(= best-performing method).%
\end{itemize}%
\end{frame}%
%
%
\section{(Multi-)Permutation Problems}%
%
\subsection{Black-Box Permutation Problems}%
%
\begin{frame}%
\frametitle{Black-Box Permutation Problems: Introduction}%
\begin{itemize}%
\item The search space~\searchSpace\ corresponds to the set of all permutations of the first \scale~natural numbers.%
%
\item<2-> The algorithms know nothing about the problem nature, they only can sample permutations~\sespel\ and compute their objective values~\objFunb{\sespel}.%
%
\item<3-> When faced with a new problem where solutions can be mapped to permutations, it makes sense to try the baseline algorithms from this field.%
%
\end{itemize}%
\end{frame}%
%
\begin{frame}%
\frametitle{Black-Box Permutation Problems: Baselines}%
\begin{itemize}%
%
\item \glsFull{algoRLS} with a unary operator that swaps two elements\uncover<2->{,}%
\item<2-> \glsShort{algoRLS} with a unary operator that swaps a binomially distributed number of elements~(similar to the \glsFull{algoEAopo} for permutations)\uncover<3->{,}%
\item<3-> \glsFull{algoSA} with the same operators and suitable exponential temperature schedule setups\uncover<4->{,}%
\item<4-> \glsFull{algoTS} with the same operator that forbids moving values to recently-used indicess\uncover<5->{, and a}%
\item<5-> \glsFull{algoEAmpl} with these operators and maybe \glsFull{algoAPcross}\cite{LKPM1997DBNTOTMGWGA,LKMID1999GAFTTSPARORAO}.%
\end{itemize}%
\end{frame}%
%
%
\subsection{Traveling Salesperson Problem}%
%
\begin{frame}%
\frametitle{Traveling Salesperson Problem: Introduction}%
\begin{itemize}%
%
\item The goal of a \glsShort{optTSP} is to find the shortest round-trip tour through \scale~cities.%
%
\item<2-> Solutions are often encoded as permutations~\sespel, where each city gets a number and the order of numbers in the permutation corresponds to the order of cities in the tour.%
%
\item<3-> The distances between the \scale~cities are given in a $\scale\times\scale$~distance matrix~\matrixDistance, where \matrixIndex{\matrixDistance}{a}{b} is the distance from city~$a$ to city~$b$.%
%
\item<4-> Hence, the objective function subject to minimization is~$\objFunb{\sespel}=\matrixIndex{\matrixDistance}{\arrayIndex{\sespel}{\scale}}{\arrayIndex{\sespel}{1}}+\sum_{i=1}^{\scale-1}\matrixIndex{\matrixDistance}{\arrayIndex{\sespel}{i}}{\arrayIndex{\sespel}{i+1}}$.%
%
\item<5-> The \glsShort{optTSP} has many applications\cite{ABCC2006TTSPACS}, including in transportation, logistics, (circuit) plotting and drilling problems\cite{GJR1991OCOPADMACS}, telescope aiming, and genome sequencing.%
%
\item<6-> The \glsShort{optTSP} is \npHard\cite{GJ1979CAIAGTTTONC,GP2002TTSPAIV,K1988TCOOP}.%
%
%
\end{itemize}%
\end{frame}%
%
\begin{frame}%
\frametitle{Traveling Salesperson Problem: Benchmarks}%
\begin{itemize}%
%
\item \gls{TSPLib} with 110~symmetric instances with scales~$\scale\in\intRange{14}{85\decSep900}$ and 19~asymmetric instances with scales~$\scale\in\intRange{17}{443}$\cite{R2016T,R1991ATSPL}.%
%
\item<2-> The instances of the 8th~DIMACS challenge\cite{JMGR20008DICTTSP,JM2002EAOHFTS,JGMYZZ2002EAOHFTS}, which have up to~10\decSep000\decSep000~cities.%
%
\item<3-> The \citetitle{C2013WT} repository\cite{C2013WT} with many real-world geographical instances, including one with all 1\decSep904\decSep711 registered populated places of the Earth at some point in time in the 2000s.%
%
\item<4-> The \gls{TSPSuite}, which offers an integrated \gls{Java} environment for benchmarking algorithms on the \gls{TSPLib} instances\cite{WCLTTCMY2014BOAAOSFFTTSP}.%
%
\end{itemize}%
\end{frame}%
%
\begin{frame}%
\frametitle{Traveling Salesperson Problem: Baselines}%
\begin{itemize}%
%
\item The \glsFull{algoLK} algorithm\cite{LK1973AEHAFTTSP}, which exchanges variable number of edges in a tour in each step.%
%
\end{itemize}%
\end{frame}%
%
\begin{frame}%
\frametitle{Traveling Salesperson Problem: SOTA}%
\begin{itemize}%
%
\item The \glsFull{algoLKH} algorithm building on top of the \glsShort{algoLK} algorithm\cite{H2009GKOSFTLKTH,H2000AEIOTLKTSH}.%
%
\item<2-> An \glsFull{algoEA} with \glsFull{algoEAX} as binary operator by \citeauthor{NK1997EACAHPGAFTTSP}~\bracketCite{NK1997EACAHPGAFTTSP,NS2012ANGAFTATSP,NK2013APGAUEACFTTSP}.%
%
\item<3-> An \glsFull{algoMA} with \glsFull{algoGPX} as binary operator by \textcite{WHH2010AHGAFTTSPUGPC}.
%
\item<4-> Especially: \glsShort{algoLKH} run in an \glsShort{algoMA}-fashion with \glsShort{algoGPX} as binary operator as proposed by \textcite{THW2018ERITLKHTSH}.%
%
\item<5-> \gls{Concorde}, an exact solver for the \gls{optTSP} by \citeauthor{ABCC2006TTSPACS}~\bracketCite{C1997CTS,ABCC2006TTSPACS,C2012IPOTTSMATLOC}, which even solved the largest instance of \gls{TSPLib}~(with 85\decSep900~cities) to optimality!\uncover<6->{ %
(This algorithm uses \glsShort{algoLK}-style algorithms methods like \glsShort{algoLKH} to find its starting point.)%
}%
%
\end{itemize}%
\end{frame}%
%
\subsection{Quadratic Assignment Problem}%
%
\begin{frame}[t]%
\frametitle{Quadratic Assignment Problem: Introduction}%
\begin{itemize}%
%
\item In a \glsFull{optQAP}, \scale~facilities need to be assigned to \scale~locations\cite{KB1957APATLOEA,BCPP1998TQAP,LdABNHQ2007ASFTQAP,CWTW2024FFAOWBFGSORLSOTQAP}.%
%
\item<2-> Such an assignment can be represented as a permutation~\sespel\ of the first \scale~natural numbers, where \arrayIndex{\sespel}{i}~specifies the location where facility~$i$ should be placed.%
%
\item<3-> \matrixDistance~is the distance matrix and \matrixIndex{\matrixDistance}{a}{b}~is the distance from location~$a$ to location~$b$.%
%
\item<4-> \matrixFlow~is the flow matrix and \matrixIndex{\matrixFlow}{i}{j}~is the amount of material flowing~(e.g., daily)~from facility~$i$ to facility~$j$.%
%
\item<5-> We want to minimize the objective function~$\objFunb{\sespel}=\sum_{i=1}^{\scale}\sum_{j=1}^{\scale} \matrixIndex{\matrixDistance}{\arrayIndex{\sespel}{i}}{\arrayIndex{\sespel}{j}}*\matrixIndex{\matrixFlow}{i}{j}$, i.e., we want to minimize the sum of the products of distance and flows.%
%
\item<6-> The \glsShort{optQAP} has applications such as building layout\cite{E1977HLAAQAP,CNTSA2021HLDRAAQAPWGD,KP1978CALD}, keyboard layout\cite{BO1977EVSMQZ}, circuit design\cite{EW1990OSOSTFSM}, wiring\cite{S1961TBWPAPA}, and scheduling\cite{S2011SISBSMSWST}.%
%
\item<7-> The \glsShort{optQAP} is \npComplete\cite{SGA1976PCAP,DPST2006MFHOSATSEAGAACMACS} and even among this problem class, is considered as particularly hard for \glspl{metaheuristic}\cite{LdABNHQ2007ASFTQAP,BAP2015MAFTQAPPAC,S2006ILSFTQAP,BCPP1998TQAP}.%
%
\end{itemize}%
%
\locateGraphic{-4}{width=0.8\paperwidth}{\sharedDir/graphics/qap_example/qap_example}{0.1}{0.57}%
%
\end{frame}%
%
\begin{frame}%
\frametitle{Quadratic Assignment Problem: Benchmarks}%
\begin{itemize}%
\item the 136~instances of \gls{QAPLib}\cite{BKR1997QAQAPL,BCKRAH2022QAQAPLPIAS}, 28~of which have not yet been solved to optimality\cite{A2023TQC}~(at least by \citeyear{A2023TQC}).%
\end{itemize}%
\end{frame}%
%
%
\begin{frame}%
\frametitle{Quadratic Assignment Problem: Baselines}%
\begin{itemize}%
\item the robust \glsFull{algoTS} algorithms by \textcite{T1991RTSFTQAP,T1995COISFTQAP} and \textcite{M2005ATSAFTQAP},%
\item<2-> \glsFull{algoSA} by \textcite{TB1994AISAAFTQAP}, \textcite{M2003AMSAAFTQAP},%
\item<3-> the \glsFull{algoRTS} by \textcite{BT1994SAATSITLRACOQT,BT1994TRTS}, and%
\item<4-> the \glsFull{algoILS} by \textcite{S2006ILSFTQAP}.%
\end{itemize}%
\end{frame}%
%
%
\begin{frame}%
\frametitle{Quadratic Assignment Problem: SOTA}%
\begin{itemize}%
\item \textcite{BAP2015MAFTQAPPAC} provide a nice experimental overview.%
\item<2-> Apart from this, we can identify the following methods\uncover<2->{:%
\begin{itemize}%
\item the \glsFullpl{algoMA} by \textcite{FF1993GHFTQAP} and \textcite{TG1997AMFTQAP}, which use the \glsShort{algoTS} from~\cite{T1991RTSFTQAP} as local search,%
\item<3-> the \gls{algoMA} by \textcite{MF1999ACOMATSAACFTQAP}, and%
\item<4-> the Iterated \glsShort{algoTS} by \textcite{M2012AIOTITSAFTQAP}.%
\end{itemize}%
}%
\end{itemize}%
\end{frame}%
%
\subsection{Job Shop Scheduling Problem}
%
\begin{frame}[t]%
\frametitle{Job Shop Scheduling Problem: Introduction}%
\begin{itemize}%
\item The \glsFull{optJSSP}~\cite{LLRKS1993SASAAC,BDP1996TJSSPCANST,CGLL1995STAIA,GLLRK1977OAAIDSASAS} is one of the most prominent and well-studied scheduling tasks. %
%
\only<-13>{%
\item<2-> A \glsShort{optJSSP} instance is defined by \alert<3>{\jsspJobs~jobs} with operations that need to be processed by each of the \alert<4>{\jsspMachines~machines} in a job-specific order.%
\item<5-> Each such operation of each job needs a specific time.%
}%
%
\only<-13>{%
\item<6-> Job~0 first goes for 10s to machine~0\only<7->{, %
then it goes for 20s to machine~1\only<8->{, %
then for 20s to machine~2\only<9->{, %
then for 40s to machine~3\only<10->{, %
and finally for 10s to machine~4}}}}.%
%
\item<11-> Job~1 first goes for 20s to machine~1\only<12->{, %
then it goes for 10s to machine~0\only<13->{{\dots} {\dots}and so on}.%
}}%
%
\item<14-> The goal is to find the assignment of jobs to machines with the smallest possible makespan, i.e., the schedule~(\gls{ganttChart}\cite{W2003GCACA,K2000SORCP}) that finishes the fastest.%
\item<15-> The \glsShort{optJSSP} is a \npHard\cite{LLRK1981RDIDSASAS,LLRKS1993SASAAC,CPW1998AROMSCAAA}.
\end{itemize}%
%
%
\locateGraphic{3}{width=0.5\paperwidth}{\sharedDir/graphics/jssp_example/jssp_example_instance_01}{0.25}{0.625}%
\locateGraphic{3}{width=0.5\paperwidth}{\sharedDir/graphics/jssp_example/jssp_example_instance_02}{0.25}{0.625}%
\locateGraphic{4}{width=0.5\paperwidth}{\sharedDir/graphics/jssp_example/jssp_example_instance_03}{0.25}{0.625}%
\locateGraphic{5,13}{width=0.5\paperwidth}{\sharedDir/graphics/jssp_example/jssp_example_instance}{0.25}{0.625}%
\locateGraphic{6}{width=0.5\paperwidth}{\sharedDir/graphics/jssp_example/jssp_example_instance_04}{0.25}{0.625}%
\locateGraphic{7}{width=0.5\paperwidth}{\sharedDir/graphics/jssp_example/jssp_example_instance_05}{0.25}{0.625}%
\locateGraphic{8}{width=0.5\paperwidth}{\sharedDir/graphics/jssp_example/jssp_example_instance_06}{0.25}{0.625}%
\locateGraphic{9}{width=0.5\paperwidth}{\sharedDir/graphics/jssp_example/jssp_example_instance_07}{0.25}{0.625}%
\locateGraphic{10}{width=0.5\paperwidth}{\sharedDir/graphics/jssp_example/jssp_example_instance_08}{0.25}{0.625}%
\locateGraphic{11}{width=0.5\paperwidth}{\sharedDir/graphics/jssp_example/jssp_example_instance_09}{0.25}{0.625}%
\locateGraphic{12}{width=0.5\paperwidth}{\sharedDir/graphics/jssp_example/jssp_example_instance_10}{0.25}{0.625}%
\locateGraphic{14-15}{width=0.7\paperwidth}{\sharedDir/graphics/jssp_example/jssp_example_gantt}{0.125}{0.43}%
\end{frame}%
%
\begin{frame}[t]%
\frametitle{Job Shop Scheduling Problem: Possible Metaheuristic Approach}%
\begin{itemize}%
\only<-10>{%
\item The vast majority of the \gls{SOTA} algorithms for the \glsShort{optJSSP} use graph-based representations for solutions\cite{BDP1996TJSSPCANST,NS1996AFTSAFTJSP}.%
\item<2-> However, schedules / \glspl{ganttChart} / graphs are complex data structures.%
\item<3-> Defining search operators for them is a bit more complicated than working on permutations or bit strings.%
\item<4-> However, it is possible to map permutations of operations to \glspl{ganttChart}.%
\item<5-> And we know how to search in a space of permutations\dots%
\item<6-> \alert{Idea} by \textcite{GTK1994SJSSPBGA}, \citeauthor{B1995AGPATJSSWGA}~\bracketCite{B1995AGPATJSSWGA,BMK1996OPRFSP}, and \textcite{SIS1997NESFSJSPBGA}:\uncover<7->{%
\begin{itemize}%
\item The search space~\searchSpace\ corresponds of permutations with repetitions~\sespel, where each of the \jsspJobs~job~IDs appears \jsspMachines~times, i.e., the problem scale is~$\scale=\jsspJobs*\jsspMachines$.%
\item<8-> The solution space~\solutionSpace\ is the set of correct \glspl{ganttChart}~\solution.%
\item<9-> A decoding function~$\decode:\searchSpace\mapsto\solutionSpace$ is designed that constructs a \gls{ganttChart}~$\solution=\decodeb{\sespel}$ from such a permutation~\sespel.%
\item<10-> Then we can use the algorithms and operators we already know\dots%
\end{itemize}%
}%
}%
\item<11-> The decoder~\decode\ processes the permutations~$\sespel\in\searchSpace$ from beginning to end and places the operations belonging to the listed jobs at their corresponding machines at the earliest possible time in a \gls{ganttChart}~$\solution\in\solutionSpace$.%
\end{itemize}%
%
\locateGraphic{12}{width=0.7\paperwidth}{\sharedDir/graphics/jssp_example/jssp_example_01}{0.15}{0.3}%
\locateGraphic{13}{width=0.7\paperwidth}{\sharedDir/graphics/jssp_example/jssp_example_02}{0.15}{0.3}%
\locateGraphic{14}{width=0.7\paperwidth}{\sharedDir/graphics/jssp_example/jssp_example_03}{0.15}{0.3}%
\locateGraphic{15}{width=0.7\paperwidth}{\sharedDir/graphics/jssp_example/jssp_example_04}{0.15}{0.3}%
\locateGraphic{16}{width=0.7\paperwidth}{\sharedDir/graphics/jssp_example/jssp_example_05}{0.15}{0.3}%
\locateGraphic{17}{width=0.7\paperwidth}{\sharedDir/graphics/jssp_example/jssp_example_06}{0.15}{0.3}%
\locateGraphic{18}{width=0.7\paperwidth}{\sharedDir/graphics/jssp_example/jssp_example_07}{0.15}{0.3}%
\locateGraphic{19}{width=0.7\paperwidth}{\sharedDir/graphics/jssp_example/jssp_example_08}{0.15}{0.3}%
\locateGraphic{20}{width=0.7\paperwidth}{\sharedDir/graphics/jssp_example/jssp_example_09}{0.15}{0.3}%
\locateGraphic{21}{width=0.7\paperwidth}{\sharedDir/graphics/jssp_example/jssp_example_10}{0.15}{0.3}%
\locateGraphic{22}{width=0.7\paperwidth}{\sharedDir/graphics/jssp_example/jssp_example_11}{0.15}{0.3}%
\locateGraphic{23}{width=0.7\paperwidth}{\sharedDir/graphics/jssp_example/jssp_example_12}{0.15}{0.3}%
\locateGraphic{24}{width=0.7\paperwidth}{\sharedDir/graphics/jssp_example/jssp_example_13}{0.15}{0.3}%
\locateGraphic{25}{width=0.7\paperwidth}{\sharedDir/graphics/jssp_example/jssp_example_14}{0.15}{0.3}%
\locateGraphic{26}{width=0.7\paperwidth}{\sharedDir/graphics/jssp_example/jssp_example_15}{0.15}{0.3}%
\locateGraphic{27}{width=0.7\paperwidth}{\sharedDir/graphics/jssp_example/jssp_example_16}{0.15}{0.3}%
\locateGraphic{28}{width=0.7\paperwidth}{\sharedDir/graphics/jssp_example/jssp_example_17}{0.15}{0.3}%
\locateGraphic{29}{width=0.7\paperwidth}{\sharedDir/graphics/jssp_example/jssp_example_18}{0.15}{0.3}%
\locateGraphic{30}{width=0.7\paperwidth}{\sharedDir/graphics/jssp_example/jssp_example_19}{0.15}{0.3}%
\locateGraphic{31}{width=0.7\paperwidth}{\sharedDir/graphics/jssp_example/jssp_example_20}{0.15}{0.3}%
\locateGraphic{32}{width=0.7\paperwidth}{\sharedDir/graphics/jssp_example/jssp_example_21}{0.15}{0.3}%
\locateGraphic{33}{width=0.7\paperwidth}{\sharedDir/graphics/jssp_example/jssp_example_22}{0.15}{0.3}%
\locateGraphic{34}{width=0.7\paperwidth}{\sharedDir/graphics/jssp_example/jssp_example_23}{0.15}{0.3}%
\locateGraphic{35}{width=0.7\paperwidth}{\sharedDir/graphics/jssp_example/jssp_example_24}{0.15}{0.3}%
%
\end{frame}%
%
%
\begin{frame}[t]%
\frametitle{Job Shop Scheduling Problem: Benchmark}%
%
\begin{itemize}%
\item The 8~most common sets of \glsShort{optJSSP} benchmark instances are~\instFormat{abz}\cite{ABZ1988TSBPFJSS}, %
\instFormat{dmu}\cite{DMU1998BFSSP}, %
\instFormat{ft}\cite{FT1963PLCOLJSSR}, %
\alert<5>{\instFormat{la}}\cite{L1984RCPSAEIOHSTS}, %
\instFormat{orb}\cite{AC1991ACSOTJSSP}, %
\alert<6>{\instFormat{swv}}\cite{SWV1992NSSFSPWATJSS}, %
\instFormat{ta}\cite{T1993BFBSP}, and %
\alert<7>{\instFormat{yn}}\cite{YN1992AGAATLSJSI}.%
%
\only<-4>{%
\item<2-> Their data can be found (sometimes partially) in several repositories in the Internet, such as\uncover<2->{%
\begin{itemize}%
\item the \gls{ORLib} managed by~\textcite{B1990OLDTPBEM,B1990WTOL} \citeurl{B1990WTOL}\uncover<3->{,}%
\item<3-> \citeauthor{S2019JSSPH}'s page \citeurl{S2019JSSPH}\uncover<4->{, or}%
\item<4-> my own repository~\cite{W2019JRDAIOTJSSP} at \citeurl{W2019JRDAIOTJSSP}, where I collect all the above problem instances and many results from existing works.%
\end{itemize}}}%
%
\end{itemize}%
%
\locateGraphic{5}{width=0.9\paperwidth}{\sharedDir/graphics/jssp_instances/jssp_instance_la24}{0.05}{0.28}%
\locateGraphic{6-}{width=0.266666667\paperwidth}{\sharedDir/graphics/jssp_instances/jssp_instance_la24}{0.05}{0.24}%
\locateGraphic{6}{width=0.583333333\paperwidth}{\sharedDir/graphics/jssp_instances/jssp_instance_swv15}{0.366666667}{0.24}%
\locateGraphic{7-}{width=0.266666667\paperwidth}{\sharedDir/graphics/jssp_instances/jssp_instance_swv15}{0.366666667}{0.28}%
\locateGraphic{7}{width=0.266666667\paperwidth}{\sharedDir/graphics/jssp_instances/jssp_instance_yn4}{0.683333334}{0.28}%
%
\end{frame}%
%
%
\begin{frame}[t]%
\frametitle{Job Shop Scheduling Problem: Baselines}%
%
\begin{itemize}%
\item See the large study by \textcite{W2019JRDAIOTJSSP} at \citeurl{W2019JRDAIOTJSSP} for old/classical methods.%
\item<2-> Good classical algorithms to outperform are the \glsFull{algoTS} algorithms by \textcite{NS1996AFTSAFTJSP} and \textcite{PM2000ATSMGBSBFTJSSP} or the \glsFull{algoLS} by \textcite{BV1998GLSWSBFJSS}.%
\end{itemize}%
\end{frame}%
%
%
\begin{frame}[t]%
\frametitle{Job Shop Scheduling Problem: SOTA}%
%
\begin{itemize}%
\item the \glsFullpl{algoMA} by \textcite{H2002PJSSP}\uncover<2->{,}%
\item<2-> the \glsFull{algoTS} algorithms by \textcite{SS2018BBOVBCACSOJSS}, \textcite{PLC2015ATSPRATSTJSSP}, and \textcite{ZSRQ2008SNROTSAATTJSSP}\uncover<3->{,}%
\item<3-> the \glsShort{algoSA}-related method by \textcite{PSV2010SJSSPUTPOBABV}\uncover<4->{,}%
\item<4-> the \glsFull{algoMA} combining continuous optimization, the permutation-based encoding, and \glsShort{algoTS} by \textcite{GR2014AEAGMWABRKGAFJSS}\uncover<5->{, and}%
\item<5-> the algorithm by \textcite{VLS2015FDSFCBS}.%
\end{itemize}%
\end{frame}%
%
\subsection{2D~Bin Packing}%
%
\begin{frame}[t]%
\frametitle{2D~Bin Packing Problem: Introduction}%
\begin{itemize}%
\item In a \glsFull{optBPP2D}, we have a set of \scale~rectangles\only<2->{ %
and we want to pack them into rectangular bins without overlap such that their edges are parallel to the bin edges\only<-3>{, with the goal to use as few bins as possible}}\cite{LMM2002TDPPAS,LMV2002RAOTDBPP,IdLMMM2021ESTFTDCAP}.%
%
\item<4-> Typical flavors of this problem include the \glsFull{optBPP2DrotFree} and the \glsFull{optBPP2DnoRotFree}.%
%
\item<5-> The \glsShort{optBPP2D} is strongly \npHard\cite{LMV2002RAOTDBPP}.%
%
\item<6-> The problem is closely related to the 2D\nobreakdashes-\gls{optCSP}, where the goal is to cut rectangles from larger rectangles\dots%
%
\item<7-> It has many practical applications, e.g., in cutting~(wood, glas)\cite{LMV2002RAOTDBPP}, packing~(transportation, warehouse)\cite{LMV2002RAOTDBPP}, newspaper page layout\cite{SH2008ALOPNP}, as well as in scheduling and build formation in additive manufacturing\cite{PSTM2024SIORAMSNASOOAMSMTT,LZ2018SBPMSWTDBPC}.%
%
\end{itemize}%
%
\locateGraphic{1}{width=0.7\paperwidth}{\sharedDir/graphics/2dbpp_example/2dbpp_example_instance_cl01-020-04_1}{0.125}{0.5}%
\locateGraphic{2}{width=0.7\paperwidth}{\sharedDir/graphics/2dbpp_example/2dbpp_example_instance_cl01-020-04_2}{0.125}{0.5}%
\locateGraphic{3}{width=0.7\paperwidth}{\sharedDir/graphics/2dbpp_example/2dbpp_example_solution_cl01-020-04}{0.125}{0.55}%
%
\end{frame}%
%
\begin{frame}[t]%
\frametitle{2D~Bin Packing Problem: Benchmarks}%
\begin{itemize}%
\item The standard benchmark set in this field is \gls{2DPackLib}\cite{IdLMM20222ATDCAPL,IdLMM20212}, which can be found at \citeurl{IdLMM20212}.\uncover<2->{%
\begin{itemize}%
\item 10~instances of type \instFormat{beng}\cite{B1982PRPAHA} with 20 and 200~items\uncover<3->{,}%
\item<3-> 10~\instFormat{class} instance sets\cite{BW1987TDFBPA,MV1998ESOTTDFBPP} with 5~group of 10~instances each, with up to 100~items\uncover<4->{, and}%
\item<4-> the \instFormat{A} instance set\cite{MAVdC2010AFMFTTDGCSP} offers 43~instances with 13~to 809~items.%
\end{itemize}}%
%
\item<5-> The \glsFull{optAsqas} problem\cite{vdBBMSB2016ASIASSTFI} asks us to put objects of sizes~$1\times2$, $2\times3$, \dots, $\scale\times(\scale+1)$ into a similarly almost-square bin without waste.\uncover<6->{ %
It has only 4~interesting instances, with bin sizes $5\times4$ for $\scale=3$, \alert<6>{$15\times16$ for $\scale=8$}, $55\times56$ for $\scale=20$, and $119\times120$ for $\scale=34$.}%
%
\end{itemize}%
\locateGraphic{6}{width=0.8\paperwidth}{\sharedDir/graphics/asqas/asqas_8}{0.1}{0.635}%
\end{frame}%
%
\begin{frame}%
\frametitle{Traveling Tournament Problem}%
\begin{itemize}%
\item \alert{Problem}:%
\item<2-> \alert{Benchmarks}:%
\item<3-> \alert{Baselines}:%
\item<4-> \alert{\gls{SOTA}}:%
\end{itemize}%
\end{frame}%
%
\begin{frame}%
\frametitle{Vehicle Routing Problems}%
\begin{itemize}%
\item \alert{Problem}:%
\item<2-> \alert{Benchmarks}:%
\item<3-> \alert{Baselines}:%
\item<4-> \alert{\gls{SOTA}}:%
\end{itemize}%
\end{frame}%
%
\section{Continuous Optimization}%
%
\begin{frame}%
\frametitle{Black-Box Continuous Optimization}%
\begin{itemize}%
\item \alert{Problem}:%
\item<2-> \alert{Benchmarks}:%
\item<3-> \alert{Baselines}:%
\item<4-> \alert{\gls{SOTA}}:%
\end{itemize}%
\end{frame}%
%
%
\section{Multi-Objective Optimization / Noisy Optimization}%
%
\begin{frame}%
\frametitle{Black-Box Multi-Objective Problems}%
\begin{itemize}%
\item \alert{Problem}:%
\item<2-> \alert{Benchmarks}:%
\item<3-> \alert{Baselines}:%
\item<4-> \alert{\gls{SOTA}}:%
\end{itemize}%
\end{frame}%
%
\begin{frame}%
\frametitle{Black-Box Noisy-Objective Problems}%
\begin{itemize}%
\item \alert{Problem}:%
\item<2-> \alert{Benchmarks}:%
\item<3-> \alert{Baselines}:%
\item<4-> \alert{\gls{SOTA}}:%
\end{itemize}%
\end{frame}%
%
%
\begin{frame}%
\frametitle{Material Flow Control for Make-to-Stock with Reorder Points}%
\begin{itemize}%
\item \alert{Problem}:%
\item<2-> \alert{Benchmarks}:%
\item<3-> \alert{Baselines}:%
\item<4-> \alert{\gls{SOTA}}:%
\end{itemize}%
\end{frame}%
%
\section{Tree-based Search Spaces}%
%
\begin{frame}%
\frametitle{Equation and Algorithm Synthesis}%
\begin{itemize}%
\item \alert{Problem}:%
\item<2-> \alert{Benchmarks}:%
\item<3-> \alert{Baselines}:%
\item<4-> \alert{\gls{SOTA}}:%
\end{itemize}%
\end{frame}%
%
\section{Summary}%
%
\begin{frame}%
\frametitle{Summary}%
\begin{itemize}%
\item Today, I gave you a brief introduction into optimization.%
\item<2-> You learned roughly what it is.%
\item<3-> Then we looked at different examples for optimization problems.%
\end{itemize}%
\end{frame}%
%
\input{\sharedDir/slides/advertisement.tex}%
%
\endPresentation%
\end{document}%%
\endinput%
%
